\section{含参积分 III}

\subsection{含参反义积分 continue}

\begin{theorem}[Weierstrass 判别法]
	若存在 $F(x)$ 使得：
	\begin{enumerate}
		\item $\abs{f(x, y)} \le F(x) \quad (x, y) \in [a, +\infty) \times [c, d]$；
		\item $\displaystyle \int_{a}^{+\infty} F(x) \dd{x}$ 收敛。
	\end{enumerate}
	那么 $\displaystyle I(y) = \int_{a}^{+\infty} f(x, y) \dd{x}$ 在 $[c, d]$ 一致收敛。
\end{theorem}

同样也有 Dirichlet 判别法和 Abel 判别法：

\begin{theorem}[Dirichlet 判别法]
	若 $f(x, y), g(x, y)$ 满足下面的条件：
	\begin{enumerate}
		\item $\displaystyle \int_{a}^{A} f(x, y) \dd{x}$ 在 $[c, d]$ 上一致有界；
		\item $g(x, y)$ 关于 $x$ 单调。
		\item $g(x, y)$ 在 $x \to +\infty$ 时在 $[c, d]$ 上一致收敛于 $0$。
	\end{enumerate}
	那么 $\displaystyle \int_{a}^{+\infty} f(x, y) g(x, y) \dd{x}$ 在 $[c, d]$ 一致收敛。
\end{theorem}

\begin{theorem}[Abel 判别法]
	若 $f(x, y), g(x, y)$ 满足下面的条件：
	\begin{enumerate}
		\item $\displaystyle \int_{a}^{A} f(x, y) \dd{x}$ 在 $[c, d]$ 上一致收敛；
		\item $g(x, y)$ 关于 $x$ 单调。
		\item $g(x, y)$ 在 $[c, d]$ 上一致有界。
	\end{enumerate}
	那么 $\displaystyle \int_{a}^{+\infty} f(x, y) g(x, y) \dd{x}$ 在 $[c, d]$ 一致收敛。
\end{theorem}

\subsection{$\Gamma$ 函数和 $B$ 函数}

\subsubsection{$\Gamma$ 函数}

\begin{definition}[$\Gamma$ 函数]
	\textbf{$\Gamma$ 函数}的定义为：
	$$
	\Gamma(s) = \int_{0}^{+\infty} x^{s - 1} \E^{-x} \dd{x} \quad (s > 0)
	$$
\end{definition}

分析一下该函数的收敛性。
\begin{itemize}
	\item $s \ge 1$ 时：$\lim\limits_{x \to \infty} \frac{x^{s - 1}}{\E^{\frac{x}{2}}} = 0$，因此 $\Gamma(s)$ 的收敛性和 $\displaystyle \int_{0}^{+\infty} \E^{-\frac{x}{2}} \dd{x}$ 相同，是收敛的。
	
	\item $0 < s < 1$ 时：
	$$
	\Gamma(s) = \int_{0}^{1} x^{s - 1} \E^{-x} \dd{x} + \int_{1}^{+\infty} x^{s - 1} \E^{-x} \dd{x}
	$$
\end{itemize}

$\Gamma$ 函数有递推式：
$$
\begin{aligned}
	\Gamma(s + 1) & = \int_{0}^{+ \infty} x^{s} \E^{-x} \dd{x} \\
	& = \ab[-\E^{-x} x^{s}]_{0}^{+\infty} + s \int_{0}^{+\infty} x^{s - 1} \E^{-x} \dd{x} \\
	& = s \Gamma(s)
\end{aligned}
$$
特别地，$\Gamma(1) = 1$，因此 $\Gamma(n + 1) = n! \ (n \in \mathbb{N})$。

\subsubsection{$B$ 函数}

\begin{definition}[$B$ 函数]
	\textbf{$B$ 函数}的定义为：
	$$
	B(p, q) = \int_{0}^{1} x^{p - 1} (1 - x)^{q - 1} \dd{x}
	$$
	它又被称为\textbf{第一类 Euler 积分}。
\end{definition}

任取 $a \in (0, 1)$，将积分化为两部分：
$$
B(p, q) = \int_{0}^{a} x^{p - 1} (1 - x)^{q - 1} \dd{x} + \int_{a}^{1} x^{p - 1} (1 - x)^{q - 1} \dd{x}
$$
那么：
\begin{itemize}
	\item $x \to 0$ 时，$x^{p - 1} (1 - x)^{q - 1} \sim x^{p - 1}$，当 $p > 0$ 时左侧积分收敛。
	\item $x \to 1$ 时，$x^{p - 1} (1 - x)^{q - 1} \sim (1 - x)^{q - 1}$，当 $q > 0$ 时右侧积分收敛。
\end{itemize}

因此 $p > 0, q > 0$ 时 $B$ 函数收敛。

$B$ 函数具有下面的性质：
\begin{itemize}
	\item 对称性：$B(p, q) = B(q, p)$。
	\item 递推关系：$\displaystyle B(p, q) = \frac{(p - 1)(q - 1)}{(p + q - 1)(p + q - 2)} B(p - 1, q - 1)$。
	\item 与 $\Gamma$ 函数的关系：$\displaystyle B(p, q) = \frac{\Gamma(p) \Gamma(q)}{\Gamma(p + q)}$。
\end{itemize}

\subsection{数值积分}
